\chapter{Identidade programável}
\noindent\textit{Vinheta.} No balcão da farmácia, a atendente aponta para um tablet: ``\emph{Olhe para a câmera, por favor.}'' Reconhecimento facial para liberar o desconto do convênio. Mais tarde, ao tentar acessar um serviço público, o aplicativo pede uma prova de vida por selfie. À noite, o banco envia um aviso: ``\emph{Por segurança, seu limite foi reduzido até nova verificação.}'' Em um só dia, sua \textbf{identidade} virou uma sequência de \emph{checks} --- e qualquer falha virou \textbf{bloqueio}.

\section{Mapa bíblico}
\begin{itemize}[leftmargin=*,noitemsep]
  \item \textbf{Selo e pertencimento} (Ef~1:13; Ap~7:2--3): Deus \emph{sela} seu povo --- identidade recebida, não negociada.
  \item \textbf{Nome e lealdade} (Ap~13:16--17; 14:1): \emph{marcas} que distinguem quem adoramos.
  \item \textbf{Justiça com o frágil} (Is~1:17; Tg~1:27): proteger vulneráveis de exclusão e injustiça sistêmica.
\end{itemize}

\section{O que é identidade programável}
Chamamos de \textbf{identidade programável} o conjunto de \emph{credenciais digitais} (CPF/IDs, senhas, biometria, \emph{tokens}, reputações) que podem ser verificadas por máquinas e conectadas a \emph{permissões} (entrar, comprar, votar, viajar) e \emph{restrições} (limites, bloqueios). É útil, mas cria \textbf{dependências} e \textbf{portas} que alguém controla.

\section{Exemplos práticos (Brasil)}
\begin{itemize}[leftmargin=*,noitemsep]
  \item \textbf{Conta gov.br}: acesso unificado a serviços públicos, com níveis de verificação que podem incluir \emph{selfie} e validações biométricas.
  \item \textbf{Documentos digitais}: CNH/CRLV digitais e outras carteiras em aplicativos oficiais; sincronização com bases do Estado.
  \item \textbf{Biometria em serviços privados}: bancos, telecom e saúde usando selfie/voz/face para liberar operações e descontos.
  \item \textbf{Pagamentos rápidos + KYC}: transferências instantâneas combinadas a cadastros \emph{conheça seu cliente} --- identidade e transação se aproximam.
\end{itemize}

\section{Exemplos práticos (internacional)}
\begin{itemize}[leftmargin=*,noitemsep]
  \item \textbf{Carteiras de identidade digitais na Europa}: iniciativas públicas para credenciais verificáveis entre países e setores.
  \item \textbf{Índia (Aadhaar)}: número de identidade com biometria, usado para benefícios e autenticações amplas.
  \item \textbf{Sistemas de reputação e pontuação}: modelos \emph{públicos e privados} que atrelam comportamento/risco a acesso a serviços.
  \item \textbf{Autenticação passkey/biométrica}: grandes plataformas movendo logins para chaves criptográficas e biometria de dispositivo.
\end{itemize}

\section{Dinâmicas de risco (o que costuma dar problema)}
\begin{itemize}[leftmargin=*,noitemsep]
  \item \textbf{Falsos negativos}: \emph{você} não passa na validação (voz rouca, luz ruim, câmera ruim) e perde acesso.
  \item \textbf{Vazamentos e correlação}: bases de dados diferentes conectadas criando perfis extensos sem seu controle.
  \item \textbf{Bloqueios automáticos}: regras antifraude que punem inocentes (conta congelada, compra negada, cadastro travado).
  \item \textbf{Exclusão silenciosa}: idosos, pessoas com deficiência, áreas sem conectividade --- quem não entra no fluxo digital fica à margem.
\end{itemize}

\begin{nicbox}
\textbf{NIC aplicado}\\
\textbf{Narrativa}: ``segurança total'' justifica coleta e verificação contínuas.\\
\textbf{Identidade}: reduzida a \emph{fatores de autenticação} (algo que você tem/é/sabe).\\
\textbf{Controle}: portas de acesso podem ser abertas/fechadas por políticas e algoritmos --- às vezes sem interlocutor humano.
\end{nicbox}

\section{Mini-exegese: o selo que liberta (Ef 1:13 vs. marcas que escravizam)}
O \textbf{selo do Espírito} é dom que \emph{confirma} nossa filiação e liberdade em Cristo. Já a \textbf{marca coercitiva} de Ap~13 liga \emph{adoração} a \emph{controle econômico}. A tecnologia pode \emph{instrumentalizar} esse controle, mas o eixo é sempre \textbf{lealdade do coração}. Por isso, formamos pessoas cuja identidade não depende do ``login'' do mundo.

\section{Ferramenta --- Protocolo pessoal de identidade (15 minutos)}
\begin{checklist}
\begin{itemize}[leftmargin=*,noitemsep]
  \item \textbf{Mapa de contas}: liste serviços críticos (banco, e-mail, nuvem, gov., saúde) e quem tem acesso de recuperação.
  \item \textbf{Autenticação com sabedoria}: use autenticação de dois fatores com app físico/OTP; evite SMS como único método.
  \item \textbf{Plano B offline}: guarde \emph{backup} de códigos de recuperação em local físico seguro; anote contatos essenciais no papel.
  \item \textbf{Higiene de dados}: revise permissões de aplicativos (câmera, microfone, localização); revogue o que não usa.
  \item \textbf{Sinais de alerta}: se pedir \emph{selfie} toda hora, questione; se negar sem recurso, documente e escale.
\end{itemize}
\end{checklist}

\section{Ferramenta --- Quando a validação falha (passo a passo)}
\begin{checklist}[title={Recuperação de acesso},breakable]
\begin{itemize}[leftmargin=*,noitemsep]
  \item \textbf{Evidência}: registre tela e horário do erro; anote protocolos.
  \item \textbf{Redundância}: tente \emph{canal alternativo} (agência física, telefone verificado, e-mail secundário).
  \item \textbf{Interlocutor humano}: peça atendimento que \emph{não dependa} da mesma biometria que falhou.
  \item \textbf{Proteção financeira}: se banco, solicite carta de \emph{não reconhecimento} e congele transações suspeitas.
  \item \textbf{Comunidade}: avise família/igreja para suporte prático (pagamentos, deslocamento, documentos).
\end{itemize}
\end{checklist}

\section{Ferramenta --- Igreja e organizações (política de identidade justa)}
\begin{checklist}[title={Princípios para líderes}]
\begin{itemize}[leftmargin=*,noitemsep]
  \item \textbf{Mínimo necessário}: colete o mínimo de dados para servir; evite armazenar documentos sensíveis.
  6
  \item \textbf{Alternativa analógica}: ofereça caminho sem smartphone/biometria para membros idosos ou vulneráveis.
  \item \textbf{Transparência}: explique por que coleta, por quanto tempo, e quem tem acesso.
  \item \textbf{Resposta a incidentes}: procedimento claro para vazamentos e perdas de acesso dos membros.
\end{itemize}
\end{checklist}

\section{Perguntas para grupos pequenos}
\begin{itemize}[leftmargin=*,noitemsep]
  \item Já ficou ``travado'' por causa de uma validação que falhou? Como resolveu?
  \item Que dados você pode \textbf{minimizar} hoje (revogar, excluir, desinstalar)?
  \item Como sua igreja pode acolher quem não consegue ``passar na máquina''?
\end{itemize}

\bigskip
\noindent\textit{Próximo capítulo: Marca, moedas e portas de compra e venda --- quando conveniência vira coerção.}
